% ============================================================================
% Section 2.6 — Simple Interest (FL/TX: financial literacy emphasis)
% CCSS 7.RP.A.3
% ============================================================================
\section{Simple Interest: Earning and Paying Interest}

\begin{stepGoal}
\begin{itemize}
    \item Use $I = Prt$ to calculate simple interest for savings and loans
    \item Compare financial options using interest calculations
\end{itemize}
\end{stepGoal}

\begin{stepVocabBox}
    \vocabItem{Principal ($P$)}{The starting amount of money saved or borrowed.}
    \vocabItem{Interest Rate ($r$)}{The annual percent charged or earned, written as a decimal.}
    \vocabItem{Simple Interest ($I$)}{Money earned on savings or owed on a loan: $I = Prt$.}
\end{stepVocabBox}

\begin{stepsCard}[How to Calculate and Compare Simple Interest]
    \stepItem{Identify $P$ (principal), $r$ (rate as a \textbf{decimal}), and $t$ (time in \textbf{years}).}
    \stepItem{Plug into the formula: $I = P \times r \times t$. Find the total: $A = P + I$.}
    \stepItem{To compare options, calculate each total and find the \textbf{difference}.}
\end{stepsCard}

% --- Worked Example 1 ---
\begin{stepExample}{Bank~A offers $3\%$ and Bank~B offers $4.5\%$. You deposit $\$1{,}000$ for $2$ years. Which earns more?}
    \stepShow{1} Both: $P = 1{,}000$, $t = 2$. Bank~A: $r = 0.03$. Bank~B: $r = 0.045$.

    \stepShow{2} Bank~A: $I = 1{,}000 \times 0.03 \times 2 = \$60$. Total: $\$1{,}060$.\\
    Bank~B: $I = 1{,}000 \times 0.045 \times 2 = \$90$. Total: $\$1{,}090$.

    \stepShow{3} Difference: $90 - 60 = \$30$.
    \stepResult{Bank~B earns $\$30$ more. A higher rate means more interest.}
\end{stepExample}

% --- Worked Example 2 ---
\begin{stepExample}{You want to earn $\$150$ in interest on $\$1{,}500$ at $2.5\%$. How many years?}
    \stepShow{1} $P = 1{,}500$, $r = 0.025$, $I = 150$, $t = ?$.

    \stepShow{2} Plug in: $150 = 1{,}500 \times 0.025 \times t$, so $150 = 37.5t$.

    \stepShow{3} Solve: $t = 150 \div 37.5 = 4$ years.
    \stepResult{It takes $4$ years to earn $\$150$.}
\end{stepExample}

\watchOut{Always convert the percent rate to a decimal first! $4\% = 0.04$, not $4$. Forgetting this gives an answer $100$ times too large.}

% ============================================================================
% PRACTICE
% ============================================================================
\newpage
\begin{stepPractice}{Simple Interest — Comparing Options}
    \resetProblems

    \stepPracticeHeader[step1Color]{Identify the Values}
    \prob You deposit $\$800$ at $5\%$ for $2$ years. How much interest do you earn? \answerBlank[2cm]
    \answer{$\$80$}

    \stepPracticeHeader[step2Color]{Use the Formula}
    \begin{multicols}{2}
    \prob You borrow $\$1{,}200$ at $6\%$ for $4$ years. How much interest do you owe? \answerBlank[2cm]
    \answer{$\$288$}
    \prob Account~A: $\$500$ at $3\%$ for $3$ years. Account~B: $\$500$ at $2\%$ for $5$ years. Which earns more interest? \answerBlank[2cm]
    \answer{B ($\$50$ vs.\ $\$45$)}
    \end{multicols}

    \stepPracticeHeader[step3Color]{Compare Financial Options}
    \prob A loan of $\$2{,}000$ charges $\$240$ interest over $3$ years. What is the rate? \answerBlank[2cm]
    \answer{$4\%$}

    \wordProblem{Mia has $\$750$ to deposit. Bank~X pays $3\%$ and Bank~Y pays $4\%$. She saves for $5$ years. How much more does Bank~Y earn?}{\$}
    \answer{$\$37.50$}
\end{stepPractice}
